\documentclass[12pt]{article}
\usepackage[utf8]{inputenc}
\usepackage[T1]{fontenc}
\usepackage{hyperref}
\title{Owenia Language Resources for K-12 Curriculum Development}
\author{Compiled by Claude Sonnet 5.5 (claude-sonnet-5-5) for Andr\'as Kornai}
\date{}
\begin{document}
\maketitle

\section{Introduction}
Owenia (ISO 639-3 wsr; Glottocode owen1244) is usually spelled Oweina in the sources found. It is also called Waisara, after the two villages where it is spoken. Wikipedia classifies it as a Kainantu language of the Eastern Highlands Province, Papua New Guinea. The search results gave conflicting speaker counts: 640 (Wikipedia, 2000), 350 (Global Recordings Network) and 1,100 for the people group (Joshua Project). Very little material exists.

\section{Primary Educational Resources}
No primers, textbooks or school programmes were found.

\section{Linguistic and Reference Materials}
\begin{itemize}
\item Wikipedia entry for Oweina: \url{https://en.wikipedia.org/wiki/Oweina_language}
\end{itemize}
No grammar or dictionary was found in the search results.

\section{Religious and Community Resources}
\begin{itemize}
\item Joshua Project profile (Waisara, Owena): \url{https://joshuaproject.net/people_groups/15827}
\end{itemize}

\section{Digital and Audio Resources}
Global Recordings Network has a language page for wsr but states that no recordings are currently available: \url{https://globalrecordings.net/en/language/wsr}

\section{Curriculum Development Context}
Joshua Project says the language is used for everyday communication, with Tok Pisin as the wider lingua franca. No documentation suitable for building a curriculum was found, so any programme would begin with community-led documentation.

\end{document}